\subsection{Conceptualizations for Design Approach}
\subsubsection{Contactless Capacitive Water Level Sensing}
A contactless sensor is a small puck that attaches to the side of the water filter. It would work by detecting when the level goes below a certain point, and sending that information to the microcontroller. It's a simple and easy way to do one level detection. It produces a binary output of either above or below the detected level.\\
\noindent\textbf{Pros}
    \begin{itemize}
    \item Low-cost
    \item No contact
    \item Small mechanical approach
    \end{itemize}
\noindent\textbf{Cons}
    \begin{itemize}
    \item Need a good environment for sensor, not humid so not great inside a fridge
    \item Users need to install device themselves, making it hard to standardize what's considered "empty"
    \item Binary output doesn't allow for percentage based information, or any form of consumption metrics. 
    \end{itemize}
    
Because of the noted cons to this design style, we decided to find another solution. The user needing to place the sensor on their own device could throw off any configuration we could complete on our side. It puts too much onto the user for it to work well. 

\subsubsection{Weight based, Load}
The design places the pitcher unto a load cell to measure the weight. We will calibrate it to the changes in weight.\\
\noindent\textbf{Pros}
    \begin{itemize}
    \item Highest accuracy of all designs
    \item Fits to many pitcher dimensions
    \end{itemize}
\noindent\textbf{Cons}
    \begin{itemize}
    \item Need a good environment for sensor, not humid so not great inside a fridge
    \end{itemize}

\subsubsection{Ultrasonic}
An ultrasonic sensor estimates the water level by measuring the echo time of sound waves reflected from the water surface. The sensor would be installed attached to the side of the filter as close to the top of the water compartment. \\
\noindent\textbf{Pros}
    \begin{itemize}
    \item Inexpensive
    \item Easy implementation
    \end{itemize}
\noindent\textbf{Cons}
    \begin{itemize}
    \item Humidity problem
    \item Mounting hardware and wiring sensor is difficult without modifications to container. 
    \end{itemize}
    
If there are modifications to the filter itself in order for a design to work, then the design could not be possible to produce. If we chose this design and needed to modify the pitcher to mount the hardware or get the wiring right, then we would have to sell pre-modified pitchers as well, otherwise users might not be able to use the device properly. Because of these reasons we decided against using this design style. 

\subsubsection{Laser}
A laser or time-of-flight sensor measures the distance from the sensor to the water surface.
The measured distance is converted into liquid height and volume. \\
\noindent\textbf{Pros}
    \begin{itemize}
    \item Continous measurements
    \item Non contact
    \end{itemize}
\noindent\textbf{Cons}
    \begin{itemize}
    \item Humidity problem
    \item Clear line of sight
    \item Mounting problems
    \end{itemize}
    
Some of the same reasons we didn't choose the ultrasonic sensor apply to this laser sensor design. On top of that in order for us to use continuous measurements, the battery power required would raise and we'd have to figure out how to power the device for long periods of time. 
